\documentclass{article}

\usepackage{amsmath, amssymb, amsthm}

\usepackage{logicpuzzle}

% YYYY-MM-DD Date format
\usepackage[style=iso]{datetime2}

% Allows use of internal commands
\makeatletter

% Place an X for an empty cell
\newcommand*\placex[2]{
    \LP@ingrid{#1}{#2}{\LP@NG@columns}{\LP@NG@rows}{nonogram}
    \LP@G@setcellcontent[hcenter,vcenter]{#1}{#2}{$\times$}
    }

% For filling a solved puzzle's empty squares with X marks
% Current solution is to fill every cell with an X and filled cells will cover it.
% Partially finished puzzles will need each desired X placed individually.
% TODO is there anything less stupid
\newcommand{\nonogramfill}[2]{
    \foreach \i in {1,...,#1}
    \foreach \j in {1,...,#2}
    {
        \placex{\i}{\j}
    }
    }


\title{Nonogram Worksheet}
\author{Michael Gadhia}
\date{\today}

\begin{document}
\maketitle

% What this document is and what it assumes
\section{Introduction}

This is an annotated worksheet of nonogram/picross puzzles, and assumes basic knowledge of how to solve one.
This will hopefully become a companion to the nonogram fact sheet once it's finished, but I have a lot of TikZ learning to do before I can make nonogram variants with the \texttt{logicpuzzle} package.
This sheet will start with a few simple nonograms and then include ones that aren't line-solvable, and hopefully one day will have color and isometric triangular puzzles (if I can typeset them in \LaTeX).

\section{Normal Puzzles}

Puzzles in this section have a single unique solution you can reach without guessing, and follow the regular rules of nonograms.

\subsection{Line-Solvability}

These are still regular puzzles, but looking at one line at a time is not enough to solve them.

\end{document}